% !TeX program = xelatex
% 本文件中的姓名、联系方式、单位、项目、数据、奖项及照片均为虚构示例。
% 请在正式使用时替换为自己的真实信息。

\documentclass{claritycv}

% 可选全局配置：\resumesetup{键=值,...}；不调用时直接使用下面的类默认值。
% color：主题色；默认 ResumeClassic；已定义颜色名 -> 改变姓名、章节线和列表标记颜色。
% font：正文风格；默认 classic；classic -> 衬线，modern -> 无衬线。
% section-font：章节字体；默认 simhei；simhei -> 传统黑体，modern -> 现代无衬线字体。
% section-font-name：自定义章节字体；默认不设置；已安装字体名 -> 覆盖 section-font。
% section-line：章节线；默认 solid；solid -> 纯色线，gradient -> 主题色渐变线。
% section-size：章节字号；默认 12.4bp；填写 LaTeX 长度 -> 调整章节标题大小。
% body-size：正文字号；默认 9.2pt；填写 LaTeX 长度 -> 调整正文大小。
% name-weight：姓名字重；默认 inherit；inherit -> 跟随章节，regular -> 常规，bold -> 加粗。
% section-weight：章节字重；默认 bold；regular -> 常规，bold -> 加粗。
% entry-style：条目样式；默认 plain；plain -> 普通行，banner -> 浅色背景和左色条。
% entry-color：条目颜色；默认 black；已定义颜色名 -> 改变条目标题和 banner 颜色。
% entry-separator：标题分隔符；默认 dash；dash -> 短横线，bar -> 竖线。
% entry-weight：条目字重；默认 bold；regular -> 常规，bold -> 加粗。
% section-before-skip：章节上方间距；默认 0.85em；LaTeX 长度越大 -> 间距越宽。
% section-line-gap：标题与章节线距离；默认 0.24em；值为 LaTeX 长度。
% section-after-skip：章节线与正文距离；默认 0.38em；值为 LaTeX 长度。
% section-line-width：章节线粗细；默认 0.55pt；值为 LaTeX 长度。
\begin{document}

% \resumeheader[键=值,...]{姓名}{联系信息}：设置头部；省略 photo 时不显示照片。
% name-size：姓名字号；默认 14.8bp；填写 LaTeX 长度 -> 仅调整本次头部的姓名大小。
% photo：照片来源；图片路径 -> 显示照片，none -> 不显示照片；只按比例缩放，不裁切透明区。
% layout：照片布局；默认 overlay；overlay -> 叠放，flow -> 参与左右分栏。
% align：flow 纵向对齐；默认 bottom；top -> 顶部，center -> 居中，bottom -> 底部。
% photo-width：照片最大宽度；默认 2.75cm；值为 LaTeX 长度。
% photo-max-height：照片最大高度；默认 3.15cm；值为 LaTeX 长度。
% photo-line-gap：照片与首个章节线距离；默认 0.18em；值为 LaTeX 长度。
% photo-x-shift：水平位移；默认 0pt；负值 -> 左，正值 -> 右。
% photo-y-shift：垂直位移；默认 0pt；负值 -> 下，正值 -> 上。
% column-gap：信息与照片栏间距；默认 1em；值为 LaTeX 长度。
% name-gap：姓名与联系信息间距；默认 0.35em；值为 LaTeX 长度。
% after-skip：头部后的间距；默认 0.15em；值为 LaTeX 长度。
\resumeheader[
  photo={images/placeholder-id-photo-transparent.png}
]{林知远}{%
  \resumecontact{\faPhone}{联系电话：138-0000-0000}
  \resumecontact{\faEnvelope}{邮箱：\href{mailto:lin.zhiyuan.ai@example.com}{\nolinkurl{lin.zhiyuan.ai@example.com}}}\\
  \resumecontact{\faGlobe}{技术博客：\href{https://example.com/lin-ai-lab}{\nolinkurl{https://example.com/lin-ai-lab}} \textcolor{ResumeAccent}{\textbf{（28万+访问量）}}}
  \\
  \resumecontact{\faGithub}{代码主页：\href{https://example.com/lin-zhiyuan-code}{\nolinkurl{https://example.com/lin-zhiyuan-code}}}\\
  \resumecontact{\faGithub}{项目贡献：\href{https://example.com/chronos-lab}{ChronosLab}\textcolor{ResumeAccent}{(\faStar\hspace{0.25em}\textbf{4.8k stars})}、\href{https://example.com/opti-grid}{OptiGrid} \textcolor{ResumeAccent}{(\faStar\hspace{0.25em}\textbf{36.2k stars})}}
}

\section{个人简介}
\begin{itemize}
  \item \textbf{个人概况：}2025年毕业·计算机与智能工程系（人工智能与数据工程专业）；AI 平台与算法工程师，曾主导/参与 8+ 个智能项目从需求调研、方案设计、模型验证到上线交付的完整流程。发表/在审智能计算方向论文 4 篇（CCF A ×1、SCI 一区 ×1、CCF B ×1、EI ×1）、开放数据挑战赛银牌、国家级学科竞赛奖项 9+ 项。
  \item \textbf{能力概况：}涵盖 Python、SQL、PyTorch、LightGBM、时序特征工程、异常检测、预测区间、约束优化、模型评测及 FastAPI / Docker 工程化部署；具备数据建模、算法设计、服务治理与业务落地能力。
\end{itemize}

\section{工作经历}

% \resumeentry[键=值,...]{主标题}{副标题}{开始时间}{结束时间}：时间连接符和间距由类生成。
% style：条目样式；默认 plain；plain -> 普通行，banner -> 带浅色背景和左色条。
% color：条目颜色；默认 black；已定义颜色名 -> 标题、左色条和 banner 背景使用该颜色。
% logo：条目图标；默认不显示；图片路径 -> 在 banner 左侧显示。
% separator：主、副标题分隔符；默认继承 entry-separator（dash）；dash -> 短横线，bar -> 竖线。
% weight：本条字重；regular -> 常规，bold -> 加粗；省略 -> 继承 entry-weight（默认 bold）。
% divider：附加分隔线；默认 none；none -> 不显示，before -> 条目前，after -> 条目后。
% date-width：右侧时间列宽度；默认 9.3em；值为 LaTeX 长度。
% text-shift：banner 文字纵向位移；默认 0.4ex；负值 -> 下，正值 -> 上。
% padding-y：banner 上下内边距；默认 0.14em；LaTeX 长度越大 -> 色块越高。
\resumeentry[style=banner,color=ResumePrimary]
  {星澜能源集团：星澜智慧能源有限公司（海岳数智技术服务有限公司）}
  {AI 应用开发工程师}
  {2025.07}{至今}
\begin{itemize}
  \item \textbf{新能源功率预测平台：}基于 Temporal Fusion Transformer 与 LightGBM 构建多尺度预测模型；覆盖 60+ 个场站，融合气象与历史功率数据，将日内预测 MAPE 由 16.4\% 降至 8.7\%，调度分析耗时降低 72\%。
  \item \textbf{智能排产优化引擎：}基于 OR-Tools 与约束建模实现订单、物料和设备协同排程，方案生成时间由 45 分钟缩短至 6 分钟，计划按期达成率由 88.5\% 提升至 96.2\%。
  \item \textbf{设备异常预警系统：}基于多变量时间序列与无监督异常检测处理日均 120 万条遥测数据，故障召回率达到 93.4\%，误报率降低 38.6\%，可提前 2--6 小时发出预警。
\end{itemize}

\resumeentry[style=banner,color=ResumeAccent]
  {云栈科技集团：云栈供应链管理有限公司（虚构数字零售企业 Top5）}
  {算法工程师（实习）}
  {2025.01}{2025.06}
\begin{itemize}
  \item 参与需求预测与库存优化模型开发，完成特征构建、滚动验证、误差分析和在线推理接口封装。
\end{itemize}

\section{项目经历}

\resumeentry[color=ResumePrimary]
  {园区能耗预测与调度平台}
  {时间序列 / 运筹优化}
  {2024.08}{2025.01}
\begin{itemize}
  \item \textbf{项目背景：}面向商业园区电、冷、热负荷构建预测与调度平台，辅助削峰填谷和设备运行计划制定。
  \item \textbf{核心技术：}Python、PyTorch、LightGBM、OR-Tools、FastAPI、Redis、Docker。
  \item \textbf{主要工作：}
    \begin{enumerate}
      \item \textbf{数据管线：}接入气象、智能电表和日历数据，完成时间对齐、缺失修复与周期特征构建。
      \item \textbf{预测模型：}组合 TFT 与梯度提升模型，引入分位数损失和滚动验证输出可信预测区间。
      \item \textbf{调度优化：}建立容量、峰谷电价和设备启停约束，将示例园区综合调度成本降低 \textcolor{ResumeAccent}{\textbf{12.4\%}}。
    \end{enumerate}
\end{itemize}

\section{个人技能}
\begin{itemize}
  \item \textbf{编程与框架：}Python、SQL、PyTorch、FastAPI、Git、Docker。
  \item \textbf{算法与优化：}时间序列预测、异常检测、特征工程、模型评测、运筹优化。
  \item \textbf{工程化能力：}FastAPI、Redis、Docker、CI/CD、监控告警与服务性能分析。
\end{itemize}

\section{教育背景}
\resumeentry
  {海棠理工大学}
  {人工智能与数据工程 \quad 本科}
  {2020.09}{2024.06}
\begin{itemize}
  \item \textbf{GPA：}3.7/4.0（专业前 10\%）。
  \item \textbf{主修课程：}数据结构、机器学习、数据库系统、自然语言处理。
\end{itemize}

\section{学科竞赛}
\begin{multicols}{2}
\begin{itemize}[label={\color{ResumePrimary}\faTrophy}]
  \item 星桥杯智能数据挑战赛全国一等奖
  \item 全国高校算法应用挑战赛二等奖
  \item 青年人工智能创新竞赛银奖
  \item 校级数学建模挑战赛一等奖
\end{itemize}
\end{multicols}

\end{document}
